\chapter{Connectivity}
%%%%%%%%%%%%%%%%%%%%%%%%%%%%%%%%%%%%%%%%%%%%%%%%%%%%%%%%%%%%%%%%%%%%%%%%%%%%%%%
\section{Connectivity}

% Generated by scripts/blueprint_restate.py from TCSlib.GraphTheory.Core.Connectivity.
% Each body is the STATEMENT only, taken from the declaration's Lean docstring:
% the one-line summary plus the verbatim book statement.  Proof, proof plan,
% significance commentary and formalisation notes are excluded by construction.

\begin{definition}[A vertex cut: a proper subset $S$ whose deletion leaves $G$ disconnected]
\lean{SimpleGraph.IsVertexCut}
\label{SimpleGraph.IsVertexCut}
\leanok
A \textbf{vertex cut}: a proper subset \texttt{S} whose deletion leaves \texttt{G}
disconnected.

A \emph{vertex cut} of $G$ is a subset $V'$ of $V$ such that $G - V'$ is disconnected. A
\emph{$k$-vertex cut} is a vertex cut of $k$ elements. A complete graph has no vertex
cut; in fact, the only graphs which do not have vertex cuts are those that contain
complete graphs as spanning subgraphs.
\end{definition}

\begin{definition}[An edge cut: a set of edges of $G$ whose deletion leaves $G$ disconnected]
\lean{SimpleGraph.IsEdgeCut}
\label{SimpleGraph.IsEdgeCut}
\leanok
An \textbf{edge cut}: a set of edges of \texttt{G} whose deletion leaves \texttt{G}
disconnected.

Recall that an edge cut of $G$ is a subset of $E$ of the form $[S, \bar{S}]$, where $S$
is a nonempty proper subset of $V$. A \emph{$k$-edge cut} is an edge cut of $k$
elements. If $G$ is nontrivial and $E'$ is an edge cut of $G$, then $G - E'$ is
disconnected.
\end{definition}

\begin{definition}[Vertex connectivity $\kappa(G)$]
\lean{SimpleGraph.vertexConnectivity}
\label{SimpleGraph.vertexConnectivity}
\leanok
\uses{SimpleGraph.IsVertexCut}
Vertex connectivity \texttt{$\kappa$(G)}. If a vertex cut exists it is the minimum cut
size; otherwise (complete graphs) it is \texttt{$\nu$ $-$ 1}.

If $G$ has at least one pair of distinct nonadjacent vertices, the \emph{connectivity}
$\kappa(G)$ of $G$ is the minimum $k$ for which $G$ has a $k$-vertex cut; otherwise, we
define $\kappa(G)$ to be $\nu - 1$. Thus $\kappa(G) = 0$ if $G$ is either trivial or
disconnected.
\end{definition}

\begin{definition}[Edge connectivity $\kappa'(G)$: the minimum size of an edge cut (\texttt{sInf $\emptyset$ = 0})]
\lean{SimpleGraph.edgeConnectivity}
\label{SimpleGraph.edgeConnectivity}
\leanok
\uses{SimpleGraph.IsEdgeCut}
Edge connectivity \texttt{$\kappa$'(G)}: the minimum size of an edge cut (\texttt{sInf
$\emptyset$ = 0}).

we then define the \emph{edge connectivity} $\kappa'(G)$ of $G$ to be the minimum $k$
for which $G$ has a $k$-edge cut. If $G$ is trivial, $\kappa'(G)$ is defined to be zero.
Thus $\kappa'(G) = 0$ if $G$ is either trivial or disconnected, and $\kappa'(G) = 1$ if
$G$ is a connected graph with a cut edge.
\end{definition}

\begin{definition}[$G$ is $k$-connected]
\lean{SimpleGraph.IsKConnected}
\label{SimpleGraph.IsKConnected}
\leanok
\uses{SimpleGraph.vertexConnectivity}
\texttt{G} is \texttt{k}-connected.

$G$ is said to be \emph{$k$-connected} if $\kappa(G) \geq k$. All nontrivial connected
graphs are 1-connected.
\end{definition}

\begin{definition}[$G$ is $k$-edge-connected]
\lean{SimpleGraph.IsKEdgeConnected}
\label{SimpleGraph.IsKEdgeConnected}
\leanok
\uses{SimpleGraph.edgeConnectivity}
\texttt{G} is \texttt{k}-edge-connected.

$G$ is said to be \emph{$k$-edge-connected} if $\kappa'(G) \geq k$. All nontrivial
connected graphs are 1-edge-connected.
\end{definition}

\begin{definition}[$v$ is a cut vertex: $G$ is connected but deleting $v$ disconnects it]
\lean{SimpleGraph.IsCutVertex}
\label{SimpleGraph.IsCutVertex}
\leanok
\texttt{v} is a \textbf{cut vertex}: \texttt{G} is connected but deleting \texttt{v} disconnects it.

A vertex $v$ of $G$ is a \emph{cut vertex} if $E$ can be partitioned into two nonempty
subsets $E_1$ and $E_2$ such that $G[E_1]$ and $G[E_2]$ have just the vertex $v$ in
common. If $G$ is loopless and nontrivial, then $v$ is a cut vertex of $G$ if and only
if $\omega(G - v) > \omega(G)$.
\end{definition}

\begin{definition}[A block: connected with no cut vertex]
\lean{SimpleGraph.IsBlock}
\label{SimpleGraph.IsBlock}
\leanok
\uses{SimpleGraph.IsCutVertex}
A \textbf{block}: connected with no cut vertex.

A connected graph that has no cut vertices is called a \emph{block}. Every block with at
least three vertices is 2-connected. A \emph{block of a graph} is a subgraph that is a
block and is maximal with respect to this property. Every graph is the union of its
blocks.
\end{definition}

\begin{definition}[Two $u$--$v$ walks are internally disjoint if they share no internal vertex]
\lean{SimpleGraph.InternallyDisjoint}
\label{SimpleGraph.InternallyDisjoint}
\leanok
Two \texttt{u}--\texttt{v} walks are \textbf{internally disjoint} if they share no
internal vertex.

A family of paths in $G$ is said to be \emph{internally-disjoint} if no vertex of
$G$ is an internal vertex of more than one path of the family.
\end{definition}

\begin{definition}[A finite family of $u$--$v$ walks, pairwise internally disjoint (family form, for Menger)]
\lean{SimpleGraph.PairwiseInternallyDisjoint}
\label{SimpleGraph.PairwiseInternallyDisjoint}
\leanok
A finite family of \texttt{u}--\texttt{v} walks, pairwise internally disjoint (family
form, for Menger).
\end{definition}

\begin{definition}[Edge subdivision: replace $e = uv$ by a length-2 path through a new vertex]
\lean{SimpleGraph.subdivide}
\label{SimpleGraph.subdivide}
\leanok
Edge subdivision: replace \texttt{e = uv} by a length-2 path through a new vertex.

It is convenient, now, to introduce the operation of subdivision of an edge. An edge $e$
is said to be \emph{subdivided} when it is deleted and replaced by a path of length two
connecting its ends, the internal vertex of this path being a new vertex.
\end{definition}

\begin{theorem}[Edge connectivity is at most the minimum degree]
\lean{SimpleGraph.edgeConnectivity_le_minDegree}
\label{SimpleGraph.edgeConnectivity_le_minDegree}
\leanok
\uses{SimpleGraph.IsEdgeCut, SimpleGraph.edgeConnectivity}
Let $G$ be a finite simple graph with at least one vertex. Then the edge connectivity of
$G$ is at most its minimum degree; that is, $\kappa'(G) \le \delta(G)$, where
$\delta(G)$ denotes the least degree of any vertex of $G$.
\end{theorem}

\begin{theorem}[Vertex connectivity is at most one less than the order]
\lean{SimpleGraph.vertexConnectivity_le_card_pred}
\label{SimpleGraph.vertexConnectivity_le_card_pred}
\leanok
\uses{SimpleGraph.IsVertexCut, SimpleGraph.vertexConnectivity}
Let $G$ be a finite simple graph on $n$ vertices. Then its vertex connectivity satisfies
\[
\kappa(G) \le n - 1.
\]
\end{theorem}

\begin{theorem}[No cut vertex forces connectivity at least two]
\lean{SimpleGraph.two_le_vertexConnectivity_of_no_cutVertex}
\label{SimpleGraph.two_le_vertexConnectivity_of_no_cutVertex}
\leanok
\uses{SimpleGraph.IsCutVertex, SimpleGraph.IsVertexCut, SimpleGraph.vertexConnectivity}
Let $G$ be a finite simple graph on a vertex set with at least three vertices. If $G$ is
connected and has no cut vertex, then its vertex connectivity satisfies $\kappa(G) \ge
2$.
\end{theorem}

\begin{theorem}[Zero edge connectivity forces zero vertex connectivity]
\lean{SimpleGraph.vertexConnectivity_eq_zero_of_edgeConnectivity_eq_zero}
\label{SimpleGraph.vertexConnectivity_eq_zero_of_edgeConnectivity_eq_zero}
\leanok
\uses{SimpleGraph.IsEdgeCut, SimpleGraph.IsVertexCut, SimpleGraph.edgeConnectivity, SimpleGraph.vertexConnectivity}
Let $G$ be a finite simple graph. If the edge connectivity of $G$ vanishes, $\kappa'(G)
= 0$, then its vertex connectivity also vanishes, $\kappa(G) = 0$.
\end{theorem}

\begin{theorem}[Lowering edge connectivity by deleting one edge]
\lean{SimpleGraph.exists_deleteEdge_edgeConnectivity_eq}
\label{SimpleGraph.exists_deleteEdge_edgeConnectivity_eq}
\leanok
\uses{SimpleGraph.IsEdgeCut, SimpleGraph.edgeConnectivity}
Let $G$ be a finite simple graph whose edge connectivity satisfies $\kappa'(G) = k+1$
for some natural number $k$. Then $G$ has an edge $e$ whose deletion yields a graph $G -
e$ with $\kappa'(G - e) = k$.
\end{theorem}

\begin{theorem}[Edge deletion drops vertex connectivity by at most one]
\lean{SimpleGraph.vertexConnectivity_le_deleteEdge_succ}
\label{SimpleGraph.vertexConnectivity_le_deleteEdge_succ}
\uses{SimpleGraph.vertexConnectivity}
Let $G$ be a finite simple graph on a vertex set $V$, and let $e$ be an unordered pair
of vertices; write $G - e$ for the graph obtained from $G$ by deleting the edge $e$
(leaving $G$ unchanged if $e$ is not an edge of $G$). Then the vertex connectivities
satisfy
\[
\kappa(G) \;\le\; \kappa(G - e) + 1 .
\]
\end{theorem}

\begin{theorem}[A bridge forces a cut vertex]
\lean{SimpleGraph.exists_isVertexCut_singleton_of_isBridge}
\label{SimpleGraph.exists_isVertexCut_singleton_of_isBridge}
\uses{SimpleGraph.IsVertexCut}
Let $G$ be a connected simple graph on a finite vertex set $V$ with at least three
vertices, and suppose the edge $\{u,v\}$ of $G$ is a bridge (that is, deleting it
disconnects $G$). Then $G$ has a cut vertex: there exists a vertex $w$ such that the
singleton $\{w\}$ is a vertex cut of $G$, meaning $\{w\}$ is a proper subset of $V$ and
the graph $G - w$ obtained by deleting $w$ is disconnected.
\end{theorem}

\begin{theorem}[Vertex connectivity is at most edge connectivity]
\lean{SimpleGraph.vertexConnectivity_le_edgeConnectivity}
\label{SimpleGraph.vertexConnectivity_le_edgeConnectivity}
\uses{SimpleGraph.edgeConnectivity, SimpleGraph.exists_deleteEdge_edgeConnectivity_eq, SimpleGraph.vertexConnectivity, SimpleGraph.vertexConnectivity_eq_zero_of_edgeConnectivity_eq_zero, SimpleGraph.vertexConnectivity_le_deleteEdge_succ}
For any finite simple graph $G$, the vertex connectivity is at most the edge
connectivity:
\[
\kappa(G) \le \kappa'(G).
\]
\end{theorem}

\begin{theorem}[Whitney's connectivity inequalities]
\lean{SimpleGraph.whitney_inequalities}
\label{SimpleGraph.whitney_inequalities}
\leanok
\uses{SimpleGraph.edgeConnectivity, SimpleGraph.edgeConnectivity_le_minDegree, SimpleGraph.vertexConnectivity, SimpleGraph.vertexConnectivity_le_edgeConnectivity}
For every finite nonempty simple graph $G$, the vertex connectivity, edge connectivity,
and minimum degree satisfy
\[
\kappa(G) \;\le\; \kappa'(G) \;\le\; \delta(G),
\]
where $\kappa(G)$ is the least size of a vertex cut (or $|V(G)| - 1$ when $G$ has no
vertex cut), $\kappa'(G)$ is the least size of an edge cut, and $\delta(G)$ is the
minimum degree of a vertex of $G$.
\end{theorem}

\begin{theorem}[Reachability in an induced subgraph from a contained walk]
\lean{SimpleGraph.reachable_induce_of_support_subset}
\label{SimpleGraph.reachable_induce_of_support_subset}
\leanok
Let $G$ be a finite simple graph on a vertex set $V$, let $s \subseteq V$, and let $u, v
\in s$. If there is a walk $p$ from $u$ to $v$ in $G$ every vertex of which lies in $s$,
then $u$ and $v$, regarded as vertices of the induced subgraph $G[s]$, are reachable
from one another in $G[s]$; that is, some walk joins them within $G[s]$.
\end{theorem}

\begin{theorem}[Two internally disjoint paths in a 2-connected graph]
\lean{SimpleGraph.exists_two_internally_disjoint_paths_of_two_connected}
\label{SimpleGraph.exists_two_internally_disjoint_paths_of_two_connected}
\uses{SimpleGraph.InternallyDisjoint, SimpleGraph.vertexConnectivity}
Let $G$ be a finite graph on at least three vertices whose connectivity satisfies
$\kappa(G) \ge 2$. Then for any two distinct vertices $u$ and $v$ there exist two
$u$–$v$ paths in $G$ that are internally disjoint and, moreover, share no edge.
\end{theorem}

\begin{theorem}[Whitney's characterisation of $2$-connected graphs]
\lean{SimpleGraph.two_connected_iff_two_internally_disjoint_paths}
\label{SimpleGraph.two_connected_iff_two_internally_disjoint_paths}
\uses{SimpleGraph.InternallyDisjoint, SimpleGraph.block_edges_on_common_cycle, SimpleGraph.exists_two_internally_disjoint_paths_of_two_connected, SimpleGraph.reachable_induce_of_support_subset, SimpleGraph.two_le_vertexConnectivity_of_no_cutVertex, SimpleGraph.vertexConnectivity}
Let $G$ be a finite simple graph on a vertex set of at least three vertices. Then the
vertex connectivity satisfies $\kappa(G) \ge 2$ if and only if for every pair of
distinct vertices $u, v$ there exist two internally disjoint $u$--$v$ paths, that is,
two paths from $u$ to $v$ whose only common vertices are the endpoints $u$ and $v$.
\end{theorem}

\begin{theorem}[Two vertices of a 2-connected graph lie on a common cycle]
\lean{SimpleGraph.two_connected_vertices_on_common_cycle}
\label{SimpleGraph.two_connected_vertices_on_common_cycle}
\uses{SimpleGraph.vertexConnectivity}
Let $G$ be a finite simple graph whose vertex connectivity satisfies $\kappa(G) \ge 2$.
Then for any two distinct vertices $u$ and $v$ of $G$ there is a cycle of $G$ whose
support contains both $u$ and $v$.
\end{theorem}

\begin{theorem}[Any two edges of a block lie on a common cycle]
\lean{SimpleGraph.block_edges_on_common_cycle}
\label{SimpleGraph.block_edges_on_common_cycle}
\uses{SimpleGraph.IsBlock}
Let $G$ be a finite block — a connected graph with no cut vertex — on at least three
vertices. Then any two distinct edges $e_1$ and $e_2$ of $G$ lie on a common cycle; that
is, there is a cycle in $G$ whose edge set contains both $e_1$ and $e_2$.
\end{theorem}

\begin{theorem}[Blocks on at least three vertices are 2-connected]
\lean{SimpleGraph.block_three_vertices_two_connected}
\label{SimpleGraph.block_three_vertices_two_connected}
\leanok
\uses{SimpleGraph.IsBlock, SimpleGraph.two_le_vertexConnectivity_of_no_cutVertex, SimpleGraph.vertexConnectivity}
Let $G$ be a block on a vertex set of at least three vertices; that is, $G$ is connected
and has no cut vertex. Then the vertex connectivity of $G$ satisfies $\kappa(G) \ge 2$.
\end{theorem}

\begin{definition}[The Harary graph $H_{m,n}$ on vertex set \texttt{ZMod n} (a circulant)]
\lean{SimpleGraph.hararyGraph}
\label{SimpleGraph.hararyGraph}
The Harary graph \texttt{H\_\{m,n\}} on vertex set \texttt{ZMod n} (a circulant).

We shall show that equality holds in (3.1) by constructing an $m$-connected graph
$H_{m,n}$ on $n$ vertices that has exactly $\{mn/2\}$ edges. The structure of $H_{m,n}$
depends on the parities of $m$ and $n$; there are three cases. \emph{Case 1} $m$ even.
Let $m = 2r$. Then $H_{2r,n}$ is constructed as follows. It has vertices $0, 1, \ldots,
n-1$ and two vertices $i$ and $j$ are joined if $i - r \leq j \leq i + r$ (where
addition is taken modulo $n$). \emph{Case 2} $m$ odd, $n$ even. Let $m = 2r + 1$. Then
$H_{2r+1,n}$ is constructed by first drawing $H_{2r,n}$ and then adding edges joining
vertex $i$ to vertex $i + (n/2)$ for $1 \leq i \leq n/2$. \emph{Case 3} $m$ odd, $n$
odd. Let $m = 2r + 1$. Then $H_{2r+1,n}$ is constructed by first drawing $H_{2r,n}$ and
then adding edges joining vertex 0 to vertices $(n-1)/2$ and $(n+1)/2$ and vertex $i$ to
vertex $i+(n+1)/2$ for $1 \le i < (n-1)/2$.
\end{definition}

\begin{theorem}[Vertex connectivity of the Harary graph]
\lean{SimpleGraph.hararyGraph_isConnectivity}
\label{SimpleGraph.hararyGraph_isConnectivity}
\uses{SimpleGraph.hararyGraph, SimpleGraph.vertexConnectivity}
Let $n \ge 1$ and let $m$ be a natural number with $m < n$. Then the Harary graph
$H_{m,n}$ on $n$ vertices has vertex connectivity $\kappa(H_{m,n}) = m$ (Harary, 1962).
\end{theorem}

\begin{theorem}[Edge count of the Harary graph]
\lean{SimpleGraph.hararyGraph_edgeCard}
\label{SimpleGraph.hararyGraph_edgeCard}
\uses{SimpleGraph.hararyGraph}
Let $m$ and $n$ be natural numbers with $n \ge 1$ and $m < n$, and let $H_{m,n}$ be the
Harary graph on the $n$ vertices $\bbz/n\bbz$. Then $H_{m,n}$ has exactly $\lceil mn/2
\rceil$ edges (equivalently, $\lfloor (mn+1)/2 \rfloor$).
\end{theorem}

\begin{theorem}[Edge count of a $k$-edge-connected graph]
\lean{SimpleGraph.edgeCard_ge_of_kEdgeConnected}
\label{SimpleGraph.edgeCard_ge_of_kEdgeConnected}
\leanok
\uses{SimpleGraph.edgeConnectivity, SimpleGraph.edgeConnectivity_le_minDegree}
Let $G$ be a finite simple graph with vertex set $V$ and edge set $E$, and let
$\kappa'(G)$ denote its edge connectivity. For every natural number $k$ with $k \le
\kappa'(G)$, \[ k\,\abs{V} \le 2\abs{E}. \]
\end{theorem}

\begin{theorem}[Connectivity equals minimum degree when the minimum degree is large]
\lean{SimpleGraph.vertexConnectivity_eq_minDegree_of_delta_ge}
\label{SimpleGraph.vertexConnectivity_eq_minDegree_of_delta_ge}
\uses{SimpleGraph.vertexConnectivity}
Let $G$ be a finite simple graph on $\nu$ vertices with minimum degree $\delta$ and
vertex connectivity $\kappa(G)$. If $\delta \ge \nu - 2$, then $\kappa(G) = \delta$.
\end{theorem}

\begin{theorem}[Equality of vertex and edge connectivity for cubic graphs]
\lean{SimpleGraph.vertexConn_eq_edgeConn_of_threeRegular}
\label{SimpleGraph.vertexConn_eq_edgeConn_of_threeRegular}
\uses{SimpleGraph.edgeConnectivity, SimpleGraph.vertexConnectivity}
Let $G$ be a finite simple graph that is $3$-regular, meaning every vertex has exactly
three neighbours. Then the vertex connectivity and the edge connectivity of $G$
coincide, $\kappa(G) = \kappa'(G)$.
\end{theorem}

\begin{theorem}[Edge connectivity at least two via edge-disjoint paths]
\lean{SimpleGraph.two_edge_connected_iff_two_edge_disjoint_paths}
\label{SimpleGraph.two_edge_connected_iff_two_edge_disjoint_paths}
\uses{SimpleGraph.edgeConnectivity}
Let $G$ be a finite simple graph with edge connectivity $\kappa'(G)$. Then $\kappa'(G)
\ge 2$ if and only if for every pair of distinct vertices $u$ and $v$ there exist two
paths from $u$ to $v$ whose edge sets are disjoint.
\end{theorem}
